\documentclass[11pt]{article}
\usepackage[margin=1in]{geometry}
\usepackage{amsmath,amssymb,amsthm}
\usepackage[T1]{fontenc}
\usepackage{lmodern}
\usepackage{microtype}
\usepackage[hidelinks]{hyperref}
\newtheorem{theorem}{Theorem}
\newcommand{\ns}{\mathbin{\#}}
\title{Cancellation for the Hessenberg natural sum\\\large Proof of conjecture 00000003522}
\author{gaochengzhecpu}
\date{October 4, 2026}
\begin{document}
\maketitle
\noindent\textbf{Result.} The natural sum is cancellative in both arguments.
Ordinary ordinal addition fails right cancellation. This is a standard fact
about ordinal arithmetic~\cite{Lipparini}, proved here for the supplied assertion.

\paragraph{Definition.}
For ordinals $\alpha,\beta$, their Hessenberg natural sum $\alpha\ns\beta$
is recursively the least ordinal strictly greater than every
$\alpha'\ns\beta$ with $\alpha'<\alpha$ and every
$\alpha\ns\beta'$ with $\beta'<\beta$. Equivalently, it is the supremum
of the successors of those ordinals, taking the supremum of the empty set
to be $0$. The operation $+$ below always denotes ordinary ordinal addition.

\begin{theorem}
For arbitrary ordinals $\alpha,\beta,\gamma$,
\begin{align*}
 \alpha\ns\beta=\alpha\ns\gamma&\quad\Longleftrightarrow\quad\beta=\gamma,\\
 \beta\ns\alpha=\gamma\ns\alpha&\quad\Longleftrightarrow\quad\beta=\gamma.
\end{align*}
Nevertheless $0+\omega=1+\omega$ while $0\ne1$.
\end{theorem}
\begin{proof}
If $\beta<\gamma$, the defining lower set of $\alpha\ns\gamma$ contains
$\alpha\ns\beta$, so $\alpha\ns\beta<\alpha\ns\gamma$.
Thus the natural sum is strictly increasing in its second argument.
Interchanging the roles of the arguments in the same definition shows
that it is strictly increasing in its first argument too.
For either asserted equality, ordinal trichotomy rules out $\beta<\gamma$
and $\gamma<\beta$, leaving $\beta=\gamma$. The converse follows by substitution.

For ordinary addition, $0+\omega=\omega$, and continuity at the limit gives
$1+\omega=\sup_{n<\omega}(1+n)=\omega$. This proves the stated failure
of right cancellation. Ordinary addition still satisfies
$\alpha+\beta=\alpha+\gamma\Rightarrow\beta=\gamma$; only the opposite
direction fails in this example.
\end{proof}

\paragraph{Lean verification and scope.}
The companion project uses Lean 4.19.0 and the exact Mathlib revision
recorded in its manifest. It uses Mathlib's actual ordinals, natural
addition, ordinary addition, and $\omega$. Strict monotonicity follows
from the library's recursive lower-set characterization; injectivity gives
both cancellation laws. The final theorem includes the explicit
ordinary-addition witness and quantifies over arbitrary ordinals in every
Lean universe, with no countability or finite-representation restriction.
The library supplies the standard recursive operation; no substitute
operation or cancellation axiom is introduced. The formal proof contains
no unfinished proof or \texttt{native\_decide}.

\begin{thebibliography}{1}\small
\bibitem{Lipparini} P.~Lipparini, \emph{An infinite natural sum},
Mathematical Logic Quarterly \textbf{62} (2016), 249--257.
See the recursive definition in Section 1 and Proposition 2.2(1).
\href{https://arxiv.org/abs/1501.05123}{arXiv:1501.05123}.
\end{thebibliography}
\end{document}
